% ========================================
% ABSTRAKT 
% ========================================

\renewcommand{\rmdefault}{ptm} % Set font to Times New Roman
\renewcommand{\sfdefault}{phv} % Set sans-serif font to Helvetica
\chapter*{\fontsize{12}{14}\selectfont Abstrakt}

\vspace{0.5cm}
\noindent
\begin{tabular}{|p{5cm}|p{10cm}|}
\hline
\textbf{Imię i nazwisko autora projektu} & Krystian Zając\\
\hline
\textbf{Rok urodzenia autora projektu} & 2003\\
\hline
\textbf{Stopień/tytuł zawodowy, imię i nazwisko prowadzącego przedmiotu} & mgr inż. Artur Pelo\\
\hline
\textbf{Instytut/Kierunek/Poziom/ Moduł specjalnościowy} & Instytut Nauk Inżynieryjno-Technicznych / Informatyka / Inżynierskie / Informatyka w przedsiębiorstwie\\
\hline
\textbf{Data realizacji projektu} & od 29.04.2026 do 10.06.2026\\
\hline
\textbf{Nazwa przedmiotu} & systemy mobilne\\
\hline
\textbf{Rok studiów/semestr} & rok III, semestr VI\\
\hline
\end{tabular}

\vspace{1.5cm}

\noindent
\begin{tabular}{|p{5cm}|p{10cm}|}
\hline
\textbf{Tytuł projektu} & Projekt aplikacji mobilnej zintegrowanej z bazą danych\\
\hline
\textbf{Słowa kluczowe (max 5)} & informatyka, aplikacje mobilne, bazy danych\\
\hline
\textbf{Streszczenie (max 1400 znaków)} & Projekt dotyczy natywnej aplikacji mobilnej na system Android, przeznaczonej do monitorowania parametrów środowiskowych, głównie temperatury. Aplikacja pobiera dane z zewnętrznego serwera przez REST API i wyświetla je w czytelnej formie. Główne funkcjonalności to: bieżący podgląd aktualnej temperatury na karcie, historia pomiarów w formie przewijalnej listy, konfiguracja adresu IP serwera, regulacja czasu odświeżania oraz wybór motywu wizualnego (jasny/ciemny/systemowy). Zastosowano nowoczesne technologie: Retrofit do obsługi komunikacji sieciowej, SharedPreferences do przechowywania ustawień użytkownika, RecyclerView do efektywnego wyświetlania list oraz Snackbar do interaktywnej obsługi błędów połączenia. Aplikacja charakteryzuje się modułową architekturą, responsywnym interfejsem zgodnym z Material Design oraz automatycznym odświeżaniem danych w konfigurowalnych odstępach czasu. System jest w pełni kompatybilny z urządzeniami z Android 8.0+ i został przetestowany pod kątem stabilności oraz wydajności. Projekt stanowi solidną bazę do dalszej rozbudowy o dodatkowe funkcjonalności, takie jak wykresy statystyczne czy system powiadomień. \\
\hline
\end{tabular}

\vspace{1.5cm}
